\chapter{Requirements and Analysis}

Chapter 2 established the motivation for this study by reviewing grammatical error correction (GEC), the limitations of accuracy-only evaluation, and the problem of unnecessary rewriting in large language model (LLM)-based systems. This chapter translates those issues into project requirements and defines the experimental design used in this dissertation.

The chapter outlines the project aim, research questions, dataset, prompt conditions, system requirements, evaluation strategy, and practical constraints. The focus is whether different prompt designs change the balance between grammatical correction performance and over-correction risk in second-language (L2) learner writing.

In this project, over-correction refers to changes made beyond what is needed for grammatical repair. A correction may therefore be grammatically acceptable but still undesirable if it unnecessarily changes the learner's wording, structure, or style without clear benefit. For this reason, the study does not treat textual change alone as over-correction. Instead, over-correction risk is considered in relation to the extent of source-output change, meaning preservation, and fluency improvement.

\section{Aims and Objectives}

The aim of this project is to examine how different prompt designs affect correction performance and over-correction risk in LLM-based grammatical error correction.

The objectives are:

\begin{itemize}
    \item to evaluate multiple prompting strategies for grammatical error correction;
    \item to compare prompt variants using both grammatical accuracy and over-correction risk indicators;
    \item to identify which prompting strategy provides the best balance between correction quality and minimal, useful editing;
    \item to examine whether sentence length affects correction behaviour;
    \item to test whether over-correction risk rankings remain stable under alternative weighting schemes.
\end{itemize}

These objectives lead to the following research questions:

\begin{itemize}
    \item \textbf{RQ1:} How do prompt variants influence grammatical correction performance and over-correction risk?
    \item \textbf{RQ2:} Which prompting strategy provides the best balance between accuracy and reduced over-correction risk?
    \item \textbf{RQ3:} How does sentence length affect correction performance and over-correction behaviour?
    \item \textbf{RQ4:} Is there a measurable trade-off between grammatical accuracy and over-correction risk?
    \item \textbf{RQ5:} Are the over-correction risk rankings stable under alternative weighting schemes?
\end{itemize}

\section{Dataset}

This project uses the W\&I + LOCNESS dataset from the BEA-2019 shared task \cite{bryant-etal-2019-bea}. The dataset is suitable because it contains learner-written English sentences with manually annotated corrections, making it compatible with standard GEC evaluation methods such as ERRANT.

The dataset is used only for evaluation. No model training or fine-tuning is performed. Each learner sentence is processed independently, allowing outputs from different prompt conditions to be compared using the same input. This controlled setup helps ensure that differences in correction behaviour can be attributed mainly to prompt design rather than to changes in data.

The dataset preparation first produced a reproducible pool of candidate sentences from the processed W\&I + LOCNESS data using random seed 42. From this pool, the main prompt-refinement and prompt-comparison experiments used a fixed 500-sentence evaluation set. This ensured that all prompt conditions were tested on the same learner sentences and reference corrections.

For the sentence-length analysis, the input sentences were divided into three groups:

\begin{itemize}
    \item \textbf{Short:} 1--10 tokens;
    \item \textbf{Medium:} 11--20 tokens;
    \item \textbf{Long:} 21 or more tokens.
\end{itemize}

A balanced subset was required so that short, medium, and long sentences could be compared fairly. Where possible, sentences were drawn from the main 500-sentence evaluation set. The final sentence-length subset contained 450 sentences, with 150 sentences in each length group. In this subset, 417 sentences overlapped with the main 500-sentence set, while 33 additional sentences were selected from the same candidate pool to complete the balanced short, medium, and long groups.

This means that the sentence-length analysis was not based on a separate dataset, but on a controlled balanced subset derived from the same prepared corpus pipeline.

Minimal preprocessing is applied so that the original surface form of each learner sentence is preserved. This is important because the study compares each corrected output with the original sentence when analysing rewriting behaviour and over-correction risk.

\section{Prompting Strategies}

Prompt design is treated as the main experimental variable. Four prompting strategies are evaluated:

\begin{enumerate}
    \item \textbf{Baseline prompt:} a simple instruction asking the model to correct grammar.
    \item \textbf{Instruction-constrained prompt:} a prompt that explicitly asks the model to correct grammatical errors only and avoid unnecessary rewriting.
    \item \textbf{Role-based prompt:} a prompt that frames the model as an English teacher helping an L2 learner.
    \item \textbf{Few-shot minimal-edit prompt:} a prompt that includes example input--output pairs demonstrating minimal grammatical correction.
\end{enumerate}

The baseline prompt provides a comparison point, while the structured prompt strategies are intended to encourage more controlled correction behaviour. Multiple variants are tested for the structured strategies so that prompt wording can be refined before the final strategy comparison.

All prompt conditions use the same dataset, model settings, and input order. This allows the experiment to isolate the effect of prompt formulation on grammatical accuracy and over-correction risk.

\section{System Requirements}

The system must support a controlled comparison of prompt outputs. The main functional requirements are:

\begin{itemize}
    \item load sentence-level learner data and reference corrections;
    \item apply each prompt condition to the same input sentences;
    \item store generated outputs with sentence IDs and prompt metadata;
    \item evaluate grammatical correction performance using ERRANT;
    \item compute over-correction risk indicators by comparing outputs with the original learner sentences;
    \item retain sentence-level outputs for later inspection and qualitative analysis.
\end{itemize}

The main non-functional requirements are:

\begin{itemize}
    \item \textbf{Reproducibility:} prompts, settings, and outputs must be stored clearly;
    \item \textbf{Consistency:} the same dataset and model parameters must be used across prompt conditions;
    \item \textbf{Traceability:} sentence-level outputs must be linked to prompt labels and evaluation results;
    \item \textbf{Ethics:} learner text must be anonymised and used only for linguistic analysis.
\end{itemize}

\section{Experimental Design}

The study follows a controlled experimental design. Learner sentences and reference corrections are first extracted from the dataset. Each sentence is then passed through every prompt condition, producing parallel corrected outputs for comparison.

The outputs are evaluated from two perspectives. The first is grammatical accuracy, where model outputs are compared against reference corrections using ERRANT. The second is over-correction risk, where each corrected output is compared with the original learner sentence using the OCI indicator implemented in Chapter 4.

The study is organised into six analyses:

\begin{enumerate}
    \item \textbf{Prompt refinement:} variants within each prompt family are compared to identify the strongest version.
    \item \textbf{Prompt strategy comparison:} the selected instruction-constrained, role-based, and few-shot prompts are compared against the baseline.
    \item \textbf{Sentence-length analysis:} results are compared across short, medium, and long sentences.
    \item \textbf{OCI robustness analysis:} alternative weighting schemes are tested to examine whether prompt rankings remain stable.
    \item \textbf{Trade-off analysis:} grammatical correction performance and over-correction risk are compared to identify aggressive, controlled, and balanced prompt behaviours.
    \item \textbf{Qualitative analysis:} selected sentence-level examples are examined to show how potential over-correction appears in practice.
\end{enumerate}

This structure allows the study to compare both prompt performance and prompt behaviour, rather than selecting a prompt only because it achieves the highest $F_{0.5}$.

\section{Evaluation Strategy}

Grammatical correction performance is evaluated using ERRANT \cite{bryant-etal-2017-automatic}. The main reported metrics are precision, recall, and $F_{0.5}$. The $F_{0.5}$ score is used because it gives greater weight to precision, which is important in GEC where unnecessary corrections can reduce output quality \cite{ng-etal-2014-conll}.

Over-correction risk is evaluated using the Composite Over-Correction Index (OCI). In this dissertation, OCI is used as an indicator of over-correction risk rather than as a direct judgement that an individual edit is unnecessary. A higher OCI indicates greater potential over-correction risk under the metric. The calculation of OCI is implemented and explained in Chapter 4.

\section{Ethical and Practical Constraints}

The project uses an existing learner corpus and does not involve collecting new participant data. Learner text is handled in anonymised form and analysed only for linguistic purposes. Since the project focuses on evaluation rather than model training, no learner data is used to fine-tune a model. The API calls use only the sentence text required for correction and do not add new participant identifiers.

A practical limitation is that the study evaluates prompt engineering within a controlled setting using one model configuration, OpenAI GPT-4o with fixed decoding settings, rather than across multiple models or datasets. Therefore, the findings should be interpreted as evidence of how prompt design affects behaviour in the evaluated setup, rather than as a universal conclusion about all LLM-based GEC systems.
